\documentclass{zhvt-classic}

% 使用本地源流明体（GenRyuMin2 TW）按文件路径加载，不依赖系统字体名
% 版框（内框／界栏／书口侧栏／鱼尾）改用与版心印章同一印泥棕色：
%   印章图片 yshs.png 实色 = RGB(135,68,52) = HTML 874434
\usepackage{xcolor}
\usepackage{tikz}
% sealbrown 双注册：xcolor 供 TikZ/pgf，L3 供 cls 版框
\definecolor{sealbrown}{HTML}{874434}
\ExplSyntaxOn
\color_set:nnn { sealbrown } { HTML } { 874434 }
\ExplSyntaxOff

\zhvtset{
  grid color       = sealbrown,
  main font        = GenRyuMin2TW-SB,
  main font path   = C:/Users/Administrator/texmf/fonts/opentype/public/GenRyuMin2TW/,
jiazhu font      = GenRyuMin2TW-M,
  jiazhu font path = C:/Users/Administrator/texmf/fonts/opentype/public/GenRyuMin2TW/,
  % 页码框容纳 4 字：正常页「九十七」+前置空 1 格=4 行恰好；
  % 100+ 页「一百零一」共 5 行超出，末字「一/二/三」被框裁掉（2026-10-04 实测），故放宽到 5
  shukou page number chars = 5,
  % 书口书名框高从默认 6 字减到 4 字，把省出的 2 字高拨给 chapter 框，
  % 使 6 字书名+卷名放进 9 字高的 chapter 框后不再压线（2026-10-05）
  shukou title chars = 4,
}

\title{史记}[陈杞世家]
\author{汉·司马迁}
\maker{闲叟漫记·重製}


% ===== 索隐述赞样式（2026-11-08 四次定稿 · 愚大人亲手改定参数）=====
% ⚠ 竖排前提：自定义字体家族必须带 RawFeature={vertical:+vert}；
%   simhei.ttf 加 +vhal 会整体加载失败静默回退，只写 +vert；simfang 同理。
\newCJKfontfamily{\shuzanfont}[Path=C:/Windows/Fonts/,RawFeature={vertical:+vert}]{simhei.ttf}
\newCJKfontfamily{\shuzanbody}[Path=C:/Windows/Fonts/,RawFeature={vertical:+vert}]{simfang.ttf}
% 题头：152x30pt 圆角 6pt 棕底白字；baseline (0,2) 使整框落入界行格子不压版框线；
% 白字节点 x 不补偿、y +8.2pt（像素级墨迹实测 dx/dy≈0，L/R 留白约 20pt）。
\newcommand{\shuzanhead}[1]{%
  \begin{tikzpicture}[baseline={(0,2)}]
    \fill[sealbrown, rounded corners=6pt]
      (-76pt,-15pt) rectangle (76pt,15pt);
    \node[text=white, inner sep=0pt, anchor=center] at (0pt,8.2pt)
      {\shuzanfont\fontsize{23}{23}\selectfont #1};
  \end{tikzpicture}%
}
% ★ 修正书口卷名位置：chapter 框起始位置改回 cls 默认值，避免全局影响。
%   只在每个卷首页通过 \juan 命令里的空命令微调该页位置（安全，不会跨页污染）。
\ExplSyntaxOn
\AtBeginDocument{
  \dim_set:Nn \l__zhvt_shukou_chapter_title_start_dim
    { \l__zhvt_shukou_title_height_dim + \l__zhvt_font_size_dim * 2 }
  \dim_set:Nn \l__zhvt_shukou_page_number_start_dim
    { \l__zhvt_shukou_chapter_title_start_dim + \l__zhvt_shukou_chapter_title_height_dim }
}
\ExplSyntaxOff

\begin{document}

% ★★ 版心（书口）排布覆写：把书名放到鱼尾之下（2026-10-04）
%
%   cls 默认版心自上而下：书名（title 静态框，在鱼尾之上）→ 鱼尾 → 篇名（chapter 动态框）→ 页码
%   本模板改为：          鱼尾（位置不动）→ 书名 → 篇名 → 页码分割线 → 页码
%
%   手法：把 cls 放在 title 静态框里的书名，改放进 chapter 动态框（鱼尾之下）；
%         title 静态框留空。书名取 \l__zhvt_title_str，即 \title{...} 的第一参数，
%         所以换书时只改 \title，不必在此重复写书名。
%
%   ⚠ 三条必须遵守的写法（2026-10-04 实测）：
%     ① 汉字必须用 \hbox{...} 包住。flowfram 会把框内容装进一个宽 14.4mm（= 一个 line_sep
%        = 书口栏宽）的纵向 minipage 并按该宽度折行，而整页靠 PDF /Rotate 90 显示、
%        框内文字本身不变向 —— 不包 \hbox 时「围炉夜话」4 字会被逐字折成 4 行，
%        在版面上就成了横排 4 字并横向溢出版框。
%     ② 竖排下**一个 \hbox 就是一整列**。书名与篇名要同处一列，必须放进同一个 \hbox；
%        分成两个 \hbox 时，后一个会被当成新的一列、甩到隔壁栏去
%        （实测「序」跑到书口栏外侧 14.4mm 的正文列上）。中间的 \quad 在框内是横向空白，
%        显示出来正好是书名与篇名之间的纵向间隔。
%     ③ ★ 内容盒要写成 **与框等宽的定宽盒** `\hbox to \l__zhvt_shukou_chapter_title_height_dim{…\hfil}`。
%        实测：chapter 框是把内容盒沿纵向**居中**摆放的，内容宽度一变，书名就跟着漂：
%        自然宽度的 \hbox（书名 38.4mm）落在 x=157.2mm，加篇名后（57.6mm）落在 x=152.4mm，
%        目录页（无篇名）与正文页的书名就差半格。做成等宽盒后，无论有没有篇名，
%        内容固定从框顶起排，实测两页均落在 x=142.8mm，不再漂。
\ExplSyntaxOn
% 把书口 chapter 框（书名+篇名）起始位置下移 1 字高，避免 4 字卷名顶到鱼尾
% ⚠ cls 的 \__zhvt_calculate_dimensions: 在 \AtBeginDocument 里执行，会重算该 dim；
%   必须把覆写挂到 begin document 之后，否则会被二次覆盖（2026-10-05 实测）。
% (旧方案已移到导言区，直接修改 chapter 框起始位置)
% 卷标题命令 #1=书口篇名(鱼尾) #2=正文卷名（含卷序，卷名内空格用 ~）
%
% ★★ 目录页两项修订（2026-10-05，愚大人指出）★★
%   ① 页码起点统一对齐：
%      原状 —— 目录项写成 \par\hspace{4em}卷名\hspace{4em}页码，页码起点 = 卷名实宽 + 4em。
%      卷名实宽随字数变：「卷一~梁惠王上」6 字 +~、「卷七~离娄上」5 字 +~，
%      差 1 个字位 → 卷七~卷十 4 条的页码整体比其余 10 条高 27.21pt（一个字），看着上下不齐。
%      改法 —— 卷名补全角空格（宽 = 1 字 = 27.215pt）补足到 **7 字符位**再写进目录，
%      14 条页码起点遂全部落在同一条竖线上（起点对齐；末字随页码字数自然参差）。
%      ⚠ 只对目录里的副本 \l_tmpc_tl 补，正文卷名仍用 \l_tmpa_tl 原样 —— 正文像素不受影响。
%   ② 卷名整体上移半字位：前导 \hspace 由 4em 改 3.5em。
%      实测原状目录块在版框内偏下（天头留白 139.6pt / 地脚留白 110.7pt），
%      上移 0.5em = 13.6pt 后上下基本均分，视觉居中。
\NewDocumentCommand { \juan } { mm } {
	\clearpage
	\tl_set:Nn \l__zhvt_chapter_title_tl { #1 }
	\tl_set:Nn \l_tmpa_tl { #2 }
	\tl_set_eq:NN \l_tmpc_tl \l_tmpa_tl
	\int_step_inline:nn { 7 - \str_count:n { #2 } }
	  { \tl_put_right:Nn \l_tmpc_tl { 　 } }
	\tl_set:Nx \l_tmpb_tl { \zhnum { page } }
	\tl_put_left:Nn \l_tmpb_tl { \hspace { 4em } }
	\tl_put_left:Nn \l_tmpb_tl { \l_tmpc_tl }
	\tl_put_left:Nn \l_tmpb_tl { \par \hspace { 3.5em } }
	\addtocontents { toc } { \l_tmpb_tl }
	\skip_horizontal:n { \l__zhvt_font_size_dim * 4 }
	\l_tmpa_tl
	\par
	\thispagestyle { empty }
}
\cs_set_protected:Npn \__zhvt_set_shukou_content: {
	% ★ 鱼尾上方留白（title 静态框）嵌入闲叟钤印：
	%   整页 /Rotate 90，故此处把图用 \rotatebox{90} 转正为竖章；
	%   尺寸卡在版心条幅内（书内约 7mm 宽 × 33.8mm 高，留余量不压鱼尾）。
	\setstaticcontents * { title } {
		\vbox_to_ht:nn { \l__zhvt_baseline_skip_dim } {
			\vskip 14.2pt
			\hbox_to_wd:nn { \l__zhvt_shukou_title_height_dim } {
				\hfil
				\rotatebox [ origin = c ] { 90 } { \includegraphics [ width = 7mm , height = 33.8mm ] { yshs.png } }
				\hfil
			}
		}
	}
% ★ 版心（书口）统一改法（2026-11-16 愚大人定）：18pt＋框内居中＋书口文字同版框色
%   ① 原 27.2pt 时 7 字长篇名（如「廉颇蔺相如列传」）末字溢出 chapter 框、压住页码；
%      改 \fontsize{18pt}{18pt}，前后 \hfil 令字串在框内居中（不随篇名字数漂移）。
%   ② \vbox{\vskip 8.2pt} 定量下移，抵消缩字号后栏宽方向 8.2pt 上偏（像素实测 y 中心 543.3 = 原版）。
%      ⚠ \raisebox / \vspace* 在 flowfram 框内无效；换字号须同步换补偿：16pt→10.1、19pt→7.4。
%   ③ chapter 与 page~number 均置 \color{sealbrown}（HTML 874434），与版框／界行同色；正文仍黑。
%   ⚠ 页码字位、钤印、鱼尾、界行、正文一律未动；五部合订本同此口径。
	\setdynamiccontents * { chapter }
	{ \color{sealbrown} \vbox { \vskip 8.2pt \hbox to \l__zhvt_shukou_chapter_title_height_dim
		{ \hfil \fontsize{18pt}{18pt}\selectfont \l__zhvt_title_str \quad \l__zhvt_chapter_title_tl \hfil } } }
	\setdynamiccontents * { page~number } { \color{sealbrown} \quad \zhnum { page } }
}
\__zhvt_set_shukou_content:
\ExplSyntaxOff

\maketitle

% 单篇，无多卷目录，不渲染目录页

\juan{陈杞世家}{卷三十六~陈杞世家第六}
\hspace*{2em}陈胡公满者。虞帝舜之后也。昔舜为庶人时。尧妻之二女。居于妫汭。其后因为氏姓。姓妫氏。舜已崩。传禹天下。而舜子商均为封国。夏后之时。或失或续。至于周武王克殷纣。乃复求舜后。得妫满。封之于陈。以奉帝舜祀。是为胡公。\\
\hspace*{2em}胡公卒。子申公犀侯立。申公卒。弟相公皋羊立。相公卒。立申公子突。是为孝公。孝公卒。子慎公圉戎立。慎公当周厉王时。慎公卒。子幽公宁立。\\
\hspace*{2em}幽公十二年。周厉王奔于彘。\\
\hspace*{2em}二十三年。幽公卒。子厘公孝立。厘公六年。周宣王即位。三十六年。厘公卒。子武公灵立。武公十五年卒。子夷公说立。是岁。周幽王即位。夷公三年卒。弟平公燮立。平公七年。周幽王为犬戎所杀。周东徙。秦始列为诸侯。\\
\hspace*{2em}二十三年。平公卒。子文公圉立。\\
\hspace*{2em}文公元年。取蔡女。生子佗。十年。文公卒。长子桓公鲍立。\\
\hspace*{2em}桓公二十三年。鲁隐公初立。二十六年。卫杀其君州吁。三十三年。鲁弑其君隐公。\\
\hspace*{2em}三十八年正月甲戌己丑。桓公鲍卒。桓公弟佗。其母蔡女。故蔡人为佗杀五父及桓公太子免而立佗。是为厉公。桓公病而乱作。国人分散。故再赴。\\
\hspace*{2em}厉公二年。生子敬仲完。周太史过陈。陈厉公使以周易筮之。卦得观之否。﹁是为观国之光。利用宾于王。此其代陈有国乎。不在此。其在异国。非此其身。在其子孙。若在异国。必姜姓。姜姓。太岳之后。物莫能两大。陈衰。此其昌乎。﹂\\
\hspace*{2em}厉公取蔡女。蔡女与蔡人乱。厉公数如蔡淫。七年。厉公所杀桓公太子免之三弟。长曰跃。中曰林。少曰杵臼。共令蔡人诱厉公以好女。与蔡人共杀厉公而立跃。是为利公。利公者。桓公子也。利公立五月卒。立中弟林。是为庄公。庄公七年卒。少弟杵臼立。是为宣公。\\
\hspace*{2em}宣公三年。楚武王卒。楚始强。十七年。周惠王娶陈女为后。\\
\hspace*{2em}二十一年。宣公后有嬖姬生子款。欲立之。乃杀其太子御寇。御寇素爱厉公子完。完惧祸及己。乃奔齐。齐桓公欲使陈完为卿。完曰。﹁羁旅之臣。幸得免负檐。君之惠也。不敢当高位。﹂桓公使为工正。齐懿仲欲妻陈敬仲。卜之。占曰。﹁是谓凤皇于飞。和鸣锵锵。有妫之后。将育于姜。五世其昌。并于正卿。八世之后。莫之与京。﹂\\
\hspace*{2em}三十七年。齐桓公伐蔡。蔡败。南侵楚。至召陵。还过陈。陈大夫辕涛涂恶其过陈。诈齐令出东道。东道恶。桓公怒。执陈辕涛涂。是岁。晋献公杀其太子申生。\\
\hspace*{2em}四十五年。宣公卒。子款立。是为穆公。穆公五年。齐桓公卒。十六年。晋文公败楚师于城濮。是岁。穆公卒。子共公朔立。共公六年。楚太子商臣弑其父成王代立。是为穆王。十一年。秦穆公卒。十八年。共公卒。子灵公平国立。\\
\hspace*{2em}灵公元年。楚庄王即位。六年。楚伐陈。十年。陈及楚平。\\
\hspace*{2em}十四年。灵公与其大夫孔宁。仪行父皆通于夏姬。衷其衣以戏于朝。泄冶谏曰。﹁君臣淫乱。民何效焉。﹂灵公以告二子。二子请杀泄冶。公弗禁。遂杀泄冶。十五年。灵公与二子饮于夏氏。公戏二子曰。﹁征舒似汝。﹂二子曰。﹁亦似公。﹂征舒怒。灵公罢酒出。征舒伏弩厩门射杀灵公。孔宁。仪行父皆奔楚。灵公太子午奔晋。征舒自立为陈侯。征舒。故陈大夫也。夏姬。御叔之妻。舒之母也。\\
\hspace*{2em}成公元年冬。楚庄王为夏征舒杀灵公。率诸侯伐陈。谓陈曰。﹁无惊。吾诛征舒而已。﹂已诛征舒。因县陈而有之。群臣毕贺。申叔时使于齐来还。独不贺。庄王问其故。对曰。﹁鄙语有之。牵牛径人田。田主夺之牛。径则有罪矣。夺之牛。不亦甚乎。今王以征舒为贼弑君。故征兵诸侯。以义伐之。已而取之。以利其地。则后何以令于天下。是以不贺。﹂庄王曰。﹁善。﹂乃迎陈灵公太子午于晋而立之。复君陈如故。是为成公。孔子读史记至楚复陈。曰。﹁贤哉楚庄王。轻千乘之国而重一言。﹂\\
\hspace*{2em}八年。楚庄王卒。二十九年。陈倍楚盟。三十年。楚共王伐陈。是岁。成公卒。子哀公弱立。楚以陈丧。罢兵去。\\
\hspace*{2em}哀公三年。楚围陈。复释之。二十八年。楚公子围弑其君郏敖自立。为灵王。\\
\hspace*{2em}三十四年。初。哀公娶郑。长姬生悼太子师。少姬生偃。二嬖妾。长妾生留。少妾生胜。留有宠哀公。哀公属之其弟司徒招。哀公病。三月。招杀悼太子。立留为太子。哀公怒。欲诛招。招发兵围守哀公。哀公自经杀。招卒立留为陈君。四月。陈使使赴楚。楚灵王闻陈乱。乃杀陈使者。使公子弃疾发兵伐陈。陈君留奔郑。九月。楚围陈。十一月。灭陈。使弃疾为陈公。\\
\hspace*{2em}招之杀悼太子也。太子之子名吴。出奔晋。晋平公问太史赵曰。﹁陈遂亡乎。﹂对曰。﹁陈。颛顼之族。陈氏得政于齐。乃卒亡。自幕至于瞽瞍。无违命。舜重之以明德。至于遂。世世守之。及胡公。周赐之姓。使祀虞帝。且盛德之后。必百世祀。虞之世未也。其在齐乎。﹂\\
\hspace*{2em}楚灵王灭陈五岁。楚公子弃疾弑灵王代立。是为平王。平王初立。欲得和诸侯。乃求故陈悼太子师之子吴。立为陈侯。是为惠公。惠公立。探续哀公卒时年而为元。空籍五岁矣。\\
\hspace*{2em}十年。陈火。十五年。吴王僚使公子光伐陈。取胡。沈而去。二十八年。吴王阖闾与子胥败楚入郢。是年。惠公卒。子怀公柳立。\\
\hspace*{2em}怀公元年。吴破楚。在郢。召陈侯。陈侯欲往。大夫曰。﹁吴新得意。楚王虽亡。与陈有故。不可倍。﹂怀公乃以疾谢吴。四年。吴复召怀公。怀公恐。如吴。吴怒其前不往。留之。因卒吴。陈乃立怀公之子越。是为愍公。\\
\hspace*{2em}愍公六年。孔子适陈。吴王夫差伐陈。取三邑而去。十三年。吴复来伐陈。陈告急楚。楚昭王来救。军于城父。吴师去。是年。楚昭王卒于城父。时孔子在陈。十五年。宋灭曹。十六年。吴王夫差伐齐。败之艾陵。使人召陈侯。陈侯恐。如吴。楚伐陈。二十一年。齐田常弑其君简公。二十三年。楚之白公胜杀令尹子西。子綦。袭惠王。叶公攻败白公。白公自杀。\\
\hspace*{2em}二十四年。楚惠王复国。以兵北伐。杀陈愍公。遂灭陈而有之。是岁。孔子卒。\\
\hspace*{2em}杞东楼公者。夏后禹之后苗裔也。殷时或封或绝。周武王克殷纣。求禹之后。得东楼公。封之于杞。以奉夏后氏祀。\\
\hspace*{2em}东楼公生西楼公。西楼公生题公。题公生谋娶公。谋娶公当周厉王时。谋娶公生武公。武公立四十七年卒。子靖公立。靖公二十三年卒。子共公立。共公八年卒。子德公立。德公十八年卒。弟桓公姑容立。桓公十七年卒。子孝公匄立。孝公十七年卒。弟文公益姑立。文公十四年卒。弟平公郁立。平公十八年卒。子悼公成立。悼公十二年卒。子隐公乞立。七月。隐公弟遂弑隐公自立。是为厘公。\\
\hspace*{2em}厘公十九年卒。子愍公维立。愍公十五年。楚惠王灭陈。十六年。愍公弟阏路弑愍公代立。是为哀公。哀公立十年卒。愍公子敕立。是为出公。出公十二年卒。子简公春立。立一年。楚惠王之四十四年。灭杞。杞后陈亡三十四年。杞小微。其事不足称述。\\
\hspace*{2em}舜之后。周武王封之陈。至楚惠王灭之。有世家言。禹之后。周武王封之杞。楚惠王灭之。有世家言。契之后为殷。殷有本纪言。殷破。周封其后于宋。齐愍王灭之。有世家言。后稷之后为周。秦昭王灭之。有本纪言。皋陶之后。或封英。六。楚穆王灭之。无谱。伯夷之后。至周武王复封于齐。曰太公望。陈氏灭之。有世家言。伯翳之后。至周平王时封为秦。项羽灭之。有本纪言。垂。益。夔。龙。其后不知所封。不见也。右十一人者。皆唐虞之际名有功德臣也。其五人之后皆至帝王。余乃为显诸侯。滕。薛。驺。夏。殷。周之间封也。小。不足齿列。弗论也。\\
\hspace*{2em}周武王时。侯伯尚千余人。及幽。厉之后。诸侯力攻相并。江。黄。胡。沈之属。不可胜数。故弗采着于传云。\\
\hspace*{2em}太史公曰。舜之德可谓至矣。禅位于夏。而后世血食者历三代。及楚灭陈。而田常得政于齐。卒为建国。百世不绝。苗裔兹兹。有土者不乏焉。至禹。于周则杞。微甚。不足数也。楚惠王灭杞。其后越王句践兴。\\

% ===== 索隐述赞（司马贞《史记索隐》）· 统一样式 =====
% 题头「索隐述赞曰」五字并入反白框；正文每四句一列，定宽盒 \hfil 居列中；
% \fontsize{24}{41} 列距对齐全书界行（121+41k）；\clearpage 使述赞独占一页。
\clearpage
\shuzanhead{索隐述赞曰}\linebreak
{\shuzanbody\fontsize{24}{41}\selectfont
\hbox to \hsize{\hfil 盛德之祀。必及百世。舜。禹馀烈。\hfil}\linebreak
\hbox to \hsize{\hfil 陈。杞是继。妫满受封。东楼纂系。\hfil}\linebreak
\hbox to \hsize{\hfil 阏路篡逆。夏姬淫嬖。二国衰微。或兴或替。\hfil}\linebreak
\hbox to \hsize{\hfil 前并後虏。皆亡楚惠。句践勃兴。田和吞噬。\hfil}\linebreak
\hbox to \hsize{\hfil 蝉联血食。岂其苗裔。\hfil}%
\par}
\end{document}
